\chapter*{不定积分每日一题}

\refer{《积分的方法与技巧》金玉明\ 顾新身\ 毛瑞庭\ /\ 编著}

\exercise{求积分$\int{\frac{\d x}{ax^2+bx+c}}$}%在花括号内输入题目

\teacher{张晓鹏}%在花括号内输入你的名字

\postdate{2023年4月3日}%在花括号内输入你想要发布那一天的日期

\result{\mbox{分情况讨论，结果如下}}%在花括号内输入本题结果

\analysis{
    关于这个的积分是经常出现的，特别是对于有理分式的积分，处理到后面，经常会遇到这类的积分，因此今天我们就讨论下这个积分的解法.

    对于分母$ax^2+bx+c$，我们可以写成形如$\qty(x+B)^2+A^2$
    \formula{ax^2+bx+c=a\qty[\qty(x+\frac{b}{2a})^2]+\frac{4ac-b^2}{4a^2}}
    或者是形如$\qty(x+B)^2-A^2$：
    \formula{ax^2+bx+c=a\qty[\qty(x+\frac{b}{2a})^2]-\frac{b^2-4ac}{4a^2}}
    下面\hypertarget{4.4}{分三种情况来处理}：

    （1）当$b^2-4ac<0$时
    \formula{
        \begin{aligned}
            \int{\frac{\d x}{ax^2+bx+c}}&=\frac{1}{a}\int{\frac{\d x}{\qty(x+\dfrac{b}{2a})^2+\dfrac{4ac-b^2}{4a^2}}}\\
            &=\frac{2}{\sqrt{4ac-b^2}}\arctan{\frac{2ax+b}{\sqrt{4ac-b^2}}}+C
        \end{aligned}
    }
    其中，主要利用的积分公式是：
    \formula{\int{\frac{1}{x^2+a^2}\d x}=\frac{1}{a}\arctan{\frac{x}{a}}+C}
    （2）当$b^2-4ac<0$时
    \formula{
        \begin{aligned}
            \int{\frac{\d x}{ax^2+bx+c}}&=\frac{1}{a}\int{\frac{\d x}{\qty(x+\dfrac{b}{2a})^2-\dfrac{b^2-4ac}{4a^2}}}\\
            &=\frac{2}{\sqrt{b^2-4ac}}\arcoth{\frac{2ax+b}{\sqrt{b^2-4ac}}}+C
        \end{aligned}
    }
    其中主要利用的积分公式是：
\formula{
    \begin{aligned}
        \int{\frac{\d x}{x^2-a^2}}&=\frac{1}{2a}\int{\qty(\frac{1}{x-a}-\frac{1}{x+a})\d x}\\
        &=\frac{1}{2a}\qty[\ln{\abs{x-a}}-\ln{\abs{x+a}}]+C\\
        &=\frac{1}{2a}\ln{\abs{\frac{x-a}{x+a}}}+C=-\frac{1}{a}\arcoth{\frac{x}{a}}+C
    \end{aligned}
}
（1）（2）两种积分主要通过配方的手段，转化为$\int{\frac{\d x}{x^2-a^2}}$和$\int{\frac{\d x}{x^2+a^2}}$的形式，其
中（2）在上面的答案里用了$\arcoth{\qty(x)}$反双曲余切函数，这个主要是为了好看，写成$\ln$的形式也是完全ok的，对于反双曲余切函数感兴趣的同学可以自行百度、知乎查一查，考试并不会有所要求.

（3）当$ax^2+bx+c$可以因式分解时
\formula{ax^2+bx+c=a\qty(x-\lambda_1)\qty(x-\lambda_2)}
那么我们就可以对原积分裂项
\formula{
    \begin{aligned}
        \int{\frac{\d x}{ax^2+bx+c}}&=\frac{1}{a}\int{\frac{\d x}{\qty(x-\lambda_1)\qty(x-\lambda_2)}}\\
        &=\frac{1}{a\qty(\lambda_1-\lambda_2)}\qty(\int{\frac{\d x}{x-\lambda_1}}-\int{\frac{\d x}{x-\lambda_2}})\\
        &=\frac{1}{a\qty(\lambda_1-\lambda_2)}\qty[\ln{\abs{x-\lambda_1}}-\ln{\abs{x-\lambda_2}}]+C\\
        &=\frac{1}{a\qty(\lambda_1-\lambda_2)}\ln{\abs{\frac{x-\lambda_1}{x-\lambda_2}}}+C
    \end{aligned}
}
到这，对于$\int{\frac{\d x}{ax^2+bx+c}}$的情况和对应求解我们都搞定了，欸，你说怎么少了一种情况的样子，如果$b^2-4ac=0$呢？其实这种情况更简单，此时
\formula{\int{\frac{\d x}{ax^2+bx+c}}=\frac{1}{a}\int{\frac{\d x}{\qty(x+\dfrac{b}{2a})^2}}}
就可以直接利用基本积分表求解啦，过程就不赘述了.

其实对于这样的积分，先看看能否因式分解，因为因式分解的计算方法运算量最小，如果发现不好分解，那么再考虑考虑配方法，转化为$\qty(x+b)^2-a^2$或$\qty(x+b)^2+a^2$的形式，分别积分即可.
}%在花括号内输入分析